En esta tarea usted deberá implementar un protocolo \textbf{enlazable} de firmas de
anillo, vale decir, un protocolo de firma de anillo donde se puede detectar si dos firmas
fueron hechas por la misma persona.\footnote{En este protocolo no se puede identificar a la persona
que hizo dos firmas de anillo. Si esto fuera posible entonces el protocolo sería rastreable.} Este protocolo
estará basado en las firmas de Schnorr y el protocolo de
firmas de anillo vistos en clases.

Desde ahora en adelante suponga que los siguientes objetos son públicos: un grupo $(G, *)$ con elemento neutro $e$,
un elemento $g \in G$, un número primo $q$ tal que $|\langle
g\rangle| = q$, una función de hash $h$ que dado un string produce un número natural, y una función de hash $h_G$ que
dado un string produce un elemento del grupo $G$. Por ejemplo, la función $h$ puede ser SHA256, interpretando el output de esta función como un número natural entre $0$ y $2^{256}-1$. De la misma forma, si consideramos el grupo $(Z_n, +)$ para
$n \geq 2^{256}$, entonces $h_G$ también podría ser la función SHA256 puesto que los elementos del grupo también son números.

Recuerde que una firma de Schnorr se define de la siguiente
forma.  La clave secreta de un
usuario $A$ es un número $x_A \in \{1, \ldots, q -1\}$ y su clave
pública es el elemento del grupo $y_A = g^{x_A}$. Si el
usuario $A$ quiere generar una firma de Schnorr para un mensaje $m$,
entonces debe realizar los siguientes pasos:
\begin{enumerate}
\item Genera al azar $r \in \{1, \ldots, q-1\}$ y calcula $c = h(g^r \| m)$ considerando
$g^r$ como un string.

\item Calcula $s = r + c \cdot x_A$.

\item Define $(c, s)$ como la firma de Schnorr de $m$.
\end{enumerate}
Un usuario $B$ puede verificar si $(c,s)$ es una firma del mensaje
$m$ hecha por el usuario $A$ chequeando que la siguiente condición se
cumpla:
\begin{align*}
c \ = \ h((g^s * y_A^{q-c}) \| m).
\end{align*}
Utilizando este protocolo, mostraremos cómo se define una firma de anillo enlazable para un
conjunto de cuatro participantes. Posteriormente, usted deberá generalizar la construcción al
caso de un número arbitrario de participantes.

Suponga que los participantes generarán firmas de anillo asociadas a un evento específico,
como una encuesta o una elección presidencial. Cada evento está identificado mediante un tag $t$,
representado por un string. Por ejemplo, el tag correspondiente a las
elecciones municipales de Chile de 2028 podría ser \texttt{"Elección municipal Chile 2028"}.
Observe que eventos distintos deben utilizar etiquetas distintas, ya que estas actúan como
identificadores únicos.

Además, suponga que se ha establecido un orden entre los participantes y que todas las firmas
de anillo se generarán respetando dicho orden. Sean $1$, $2$, $3$ y $4$ los participantes que
generarán firmas de anillo. La clave secreta del usuario $i$ es $x_i$, mientras que su clave
pública es $y_i = g^{x_i}$. Finalmente, suponiendo que $r = \frac{|G|}{q}$ y considerando
cada clave pública $y_i$ como un string, defina:
\begin{align*}
T \ = \ h_G(y_1 \| y_2 \| y_3 \| y_4 \| t)^r.
\end{align*}
Dada la definición de $h_G$, se tiene que $T \in G$. Si $T = e$ (el elemento neutro del grupo), entonces
se debe utilizar un tag distinto para el evento hasta lograr que $T \neq e$. Observe que esto asegura que
$|\langle T \rangle| = q$.

Una vez definidos el tag $t$, las claves secretas y públicas de los participantes, y el elemento $T$ del grupo,
las firmas de anillo para este evento se generan de la siguiente forma.
Si el usuario $2$ quiere generar una firma de anillo de un mensaje $m$, define $f = T^{x_2}$ y realiza
los siguientes pasos:
\begin{enumerate}
\item Genera al azar $r_2 \in \{1, \ldots, q-1\}$ y calcula $c_3 = h(g^{r_2} \| m \| T^{r_2} )$.

\item Genera al azar $s_3 \in \{1, \ldots, q-1\}$.

\item Calcula $c_4 = h((g^{s_3} * y_3^{q - c_3}) \| m \| (T^{s_3} * f^{q - c_3}))$.

\item Genera al azar $s_4 \in \{1, \ldots, q-1\}$.

\item Calcula $c_1 = h((g^{s_4} * y_4^{q - c_4}) \| m \| (T^{s_4} * f^{q - c_4}))$.

\item Genera al azar $s_1 \in \{1, \ldots, q-1\}$.

\item Calcula $c_2 = h((g^{s_1} * y_1^{q - c_1}) \| m \| (T^{s_1} * f^{q - c_1}))$.

\item Calcula $s_2 = r_2 + c_2 \cdot x_2$. Note que este paso lo puede realizar el usuario $2$ ya que
conoce la clave secreta $x_2$.

\item Define $(s_1, s_2, s_3, s_4, c_1, f)$ como la firma de anillo de $m$. En particular, siempre se incluye $c_1$ en la
firma, independientemente de quién la haya generado.
\end{enumerate}
Note que cuando un participante genera una firma de anillo debe respetar el orden preestablecido. Por ejemplo, si el
usuario 3 quiere firmar, entonces debe realizar el protocolo en el orden $3 \to 4 \to 1 \to 2 \to 3$.
Además, note que como fue mencionado en clases, basta que la firma incluya los valores $s_1$, $s_2$, $s_3$, $s_4$
y un valor $c_i$ para verificar que es correcta. En este caso, este valor siempre será $c_1$, de manera tal de no
entregar información sobre quién generó la firma. Finalmente, observe que se incluye en la firma un valor $f$ adicional
que será utilizado para la enlazabilidad.

Si un usuario $B$ quiere verificar que $(s_1, s_2, s_3, s_4, c_1, f)$ es una firma
de anillo válida para el tag $t$ y el grupo formado por los usuarios $1$, $2$, $3$ y $4$, entonces debe hacer lo siguiente:
\begin{enumerate}
	\item Calcula $T \ = \ h_G(y_1 \| y_2 \| y_3 \| y_4 \| t)^r$, donde $r = \frac{|G|}{q}$
	\item Calcula $c_2=h((g^{s_1} * y_1^{q - c_1}) \| m \| (T^{s_1} * f^{q - c_1}))$
	\item Calcula $c_3=h((g^{s_2} * y_2^{q - c_2}) \| m \| (T^{s_2} * f^{q - c_2}))$
	\item Calcula $c_4=h((g^{s_3} * y_3^{q - c_3}) \| m \| (T^{s_3} * f^{q - c_3}))$
	\item Verifica si $h((g^{s_4} * y_4^{q - c_4}) \| m \| (T^{s_4} * f^{q - c_4}))$
	es igual al valor $c_1$ de la firma.
\end{enumerate}
Si esta última condición es cierta, entonces $B$ sabe
que la firma es válida, pero no sabe cuál de los usuarios firmó el
mensaje $m$.

Además, $B$ puede verificar si dos firmas fueron hechas por el mismo participante de la siguiente forma.
Si un usuario $i$ firma dos mensajes $m_1$ y $m_2$, entonces la firma de $m_1$
es de la forma $(s_1, s_2, s_3, s_4, c_1, f)$ y la firma de $m_2$ es de la forma
$(s'_1, s'_2, s'_3, s'_4, c'_1, f)$, donde en ambos casos $f =  T^{x_i}$. De hecho, todas las firmas de anillo
del usuario $i$ para este evento y este grupo van a tener como último componente $f = T^{x_i}$. De esta forma, $B$ puede
verificar si dos firmas fueron hechas por el mismo participante simplemente chequeando que las últimas
componentes en estas firmas son iguales.

Finalmente, observe que el protocolo anterior sólo nos permite verificar si dos firmas de anillo fueron generadas
por el mismo usuario en el mismo evento. En particular,
un participante puede generar firmas de anillo para eventos distintos sin que ellas sean identificadas como hechas
por la misma persona. En efecto, si consideramos dos eventos
distintos con tags $t_1 \neq t_2$, entonces esperamos que $T_1 \neq T_2$ puesto que estamos usando
una función de hash $h_G$ que es resistente a colisiones. De esta forma, una firma de anillo de un participante $i$ en
el primer evento es de la forma $(s_1, s_2, s_3, s_4, c_1, f_1)$ con $f_1 = T_1^{x_i}$, mientras que una firma
de anillo en el segundo evento es de la forma $(s'_1, s'_2, s'_3, s'_4, c'_1, f_2)$ con $f_2 = T_2^{x_i}$.
Como $f_1 \neq f_2$, puesto que $T_1 \neq T_2$, estas firmas no pueden ser reconocidas como hechas por
el mismo participante.
